\documentclass[aspectratio=169,11pt]{beamer}

\usetheme{Madrid}
\usefonttheme{professionalfonts}
\usepackage[T1]{fontenc}
\usepackage[utf8]{inputenc}
\usepackage[french]{babel}
\usepackage{lmodern}
\usepackage{tikz}
\usepackage{listings}
\usepackage{xcolor}
\usepackage{array}

\definecolor{SenegalGreen}{HTML}{007A3D}
\definecolor{SenegalYellow}{HTML}{FCD116}
\definecolor{SenegalRed}{HTML}{CE1126}
\definecolor{Ink}{HTML}{10201A}
\definecolor{SoftGreen}{HTML}{EAF7EF}
\definecolor{Panel}{HTML}{F7FAF8}
\definecolor{CodeBg}{HTML}{101820}

\setbeamercolor{normal text}{fg=Ink,bg=white}
\setbeamercolor{structure}{fg=SenegalGreen}
\setbeamercolor{title}{fg=white,bg=SenegalGreen}
\setbeamercolor{frametitle}{fg=white,bg=SenegalGreen}
\setbeamercolor{block title}{fg=white,bg=SenegalGreen}
\setbeamercolor{block body}{fg=Ink,bg=SoftGreen}
\setbeamercolor{alerted text}{fg=SenegalRed}
\setbeamertemplate{navigation symbols}{}
\setbeamertemplate{itemize item}{\color{SenegalGreen}$\blacktriangleright$}
\setbeamertemplate{itemize subitem}{\color{SenegalYellow}$\bullet$}

\setbeamertemplate{footline}{
  \leavevmode
  \hbox{
    \begin{beamercolorbox}[wd=.72\paperwidth,ht=2.4ex,dp=1ex,leftskip=1.2em]{author in head/foot}
      \color{Ink} Systeme de paiement BasicCard en FCFA
    \end{beamercolorbox}
    \begin{beamercolorbox}[wd=.28\paperwidth,ht=2.4ex,dp=1ex,rightskip=1.2em plus1fil]{date in head/foot}
      \hfill \color{Ink}\insertframenumber/\inserttotalframenumber
    \end{beamercolorbox}
  }
}

\lstdefinelanguage{BasicCard}{
  morekeywords={Option,Explicit,Const,Eeprom,As,Long,Byte,Command,End,If,Then,ElseIf,Else,Exit,Rem,Declare,Call,Sub,While,Wend,Private,Print,Line,Input},
  sensitive=false,
  morecomment=[l]{'},
  morecomment=[l]{Rem},
  morestring=[b]"
}

\lstset{
  language=BasicCard,
  basicstyle=\ttfamily\scriptsize\color{white},
  keywordstyle=\color{SenegalYellow}\bfseries,
  commentstyle=\color{gray!35},
  stringstyle=\color{green!45!white},
  backgroundcolor=\color{CodeBg},
  breaklines=true,
  frame=single,
  rulecolor=\color{SenegalGreen},
  tabsize=2,
  showstringspaces=false
}

\newcommand{\fcfa}{\textsc{FCFA}}
\newcommand{\tagbox}[2]{
  \begin{tikzpicture}
    \node[rounded corners=4pt, fill=#1, text=white, inner xsep=8pt, inner ysep=5pt] {\textbf{#2}};
  \end{tikzpicture}
}

\title[BasicCard FCFA]{Systeme de paiement BasicCard en \fcfa}
\subtitle{Carte a puce, terminal de paiement et simulation pour le Senegal}
\author{Projet BasicCard}
\institute{Simulation de paiement electronique}
\date{\today}

\begin{document}

% --------------------------------------------------------------------
\begin{frame}[plain]
  \begin{tikzpicture}[remember picture,overlay]
    \fill[SenegalGreen] (current page.south west) rectangle (current page.north east);
    \fill[SenegalYellow] (current page.south west) rectangle ([xshift=.18\paperwidth]current page.north west);
    \fill[SenegalRed] ([xshift=.82\paperwidth]current page.south west) rectangle (current page.north east);
    \node[align=center,text=white] at (current page.center) {
      {\Huge\bfseries Systeme de paiement}\\[0.35cm]
      {\Huge\bfseries BasicCard en FCFA}\\[0.45cm]
      {\large Carte a puce programmable et terminal de paiement}\\[0.75cm]
      \tagbox{Ink}{Senegal} \hspace{0.2cm}
      \tagbox{Ink}{PIN securise} \hspace{0.2cm}
      \tagbox{Ink}{Ticket de paiement}
    };
  \end{tikzpicture}
\end{frame}

% --------------------------------------------------------------------
\begin{frame}{Contexte}
  \begin{columns}[T]
    \begin{column}{0.54\textwidth}
      \begin{block}{Idee du projet}
        Simuler un paiement electronique avec une carte a puce programmable BasicCard et un terminal de paiement.
      \end{block}
      \vspace{0.2cm}
      \begin{itemize}
        \item Monnaie utilisee : francs CFA.
        \item Zone cible : Senegal.
        \item Test sans carte physique grace au simulateur BasicCard.
      \end{itemize}
    \end{column}
    \begin{column}{0.42\textwidth}
      \begin{tikzpicture}[scale=0.95]
        \fill[Panel] (0,0) rectangle (5.2,3.4);
        \draw[SenegalGreen, line width=1.2pt, rounded corners=8pt] (.4,.55) rectangle (4.8,2.85);
        \draw[SenegalYellow, line width=1.2pt] (1,.95) rectangle (2.05,1.65);
        \node[align=left] at (3.15,1.65) {\textbf{Carte FCFA}\\PIN 1234\\Solde 10000};
        \node[SenegalGreen] at (2.6,.18) {\small Carte a puce simulee};
      \end{tikzpicture}
    \end{column}
  \end{columns}
\end{frame}

% --------------------------------------------------------------------
\begin{frame}{Objectifs}
  \begin{columns}[T]
    \begin{column}{0.48\textwidth}
      \begin{block}{Objectifs techniques}
        \begin{itemize}
          \item Programmer une carte BasicCard.
          \item Declarer des commandes APDU.
          \item Stocker un solde en EEPROM.
          \item Compiler et tester dans ZCMSim.
        \end{itemize}
      \end{block}
    \end{column}
    \begin{column}{0.48\textwidth}
      \begin{block}{Objectifs fonctionnels}
        \begin{itemize}
          \item Verifier le code PIN.
          \item Consulter le solde.
          \item Payer en FCFA.
          \item Recharger la carte.
          \item Afficher un ticket de paiement.
        \end{itemize}
      \end{block}
    \end{column}
  \end{columns}
\end{frame}

% --------------------------------------------------------------------
\begin{frame}{Architecture}
  \centering
  \begin{tikzpicture}[node distance=2.2cm,>=stealth,thick]
    \tikzstyle{box}=[rounded corners=7pt, draw=SenegalGreen, fill=SoftGreen, minimum width=3.2cm, minimum height=1.15cm, align=center]
    \node[box] (user) {Utilisateur\\PIN + montant};
    \node[box, right of=user, xshift=2.5cm] (term) {Terminal\\MonTerminal.bas};
    \node[box, right of=term, xshift=2.5cm] (card) {Carte\\MaCarte.bas};
    \node[box, below of=term, yshift=-.25cm] (sim) {ZCMSim\\lecteur virtuel 201};

    \draw[->,SenegalGreen] (user) -- node[above]{saisie} (term);
    \draw[->,SenegalGreen] (term) -- node[above]{APDU} (card);
    \draw[->,SenegalGreen] (card) -- node[below]{SW1SW2 + solde} (term);
    \draw[->,SenegalRed] (sim) -- (term);
    \draw[->,SenegalRed] (sim) -- (card);
  \end{tikzpicture}
\end{frame}

% --------------------------------------------------------------------
\begin{frame}{Structure du projet}
  \begin{block}{Fichiers principaux}
    \begin{tabular}{>{\bfseries}p{0.35\textwidth} p{0.55\textwidth}}
      src/card/MaCarte.bas & Code embarque dans la carte. \\
      src/terminal/MonTerminal.bas & Interface terminal et menu utilisateur. \\
      compiler.bat & Compilation carte, terminal exe et terminal img. \\
      lancer.bat & Lancement du simulateur avec carte et terminal. \\
      docs/protocole.md & Description du protocole APDU. \\
      docs/explication-projet.md & Explication complete pour la soutenance. \\
    \end{tabular}
  \end{block}
\end{frame}

% --------------------------------------------------------------------
\begin{frame}{Fonctionnalites}
  \begin{columns}[T]
    \begin{column}{0.32\textwidth}
      \begin{block}{Paiement}
        Debit d'un montant en \fcfa{} si le solde est suffisant.
      \end{block}
    \end{column}
    \begin{column}{0.32\textwidth}
      \begin{block}{Rechargement}
        Credit d'un montant positif sur la carte.
      \end{block}
    \end{column}
    \begin{column}{0.32\textwidth}
      \begin{block}{Securite}
        PIN obligatoire et blocage apres 3 erreurs.
      \end{block}
    \end{column}
  \end{columns}
  \vspace{0.45cm}
  \centering
  \tagbox{SenegalGreen}{Solde initial : 10000 FCFA}
  \hspace{0.3cm}
  \tagbox{SenegalRed}{PIN par defaut : 1234}
\end{frame}

% --------------------------------------------------------------------
\begin{frame}{Protocole APDU}
  \begin{center}
    \renewcommand{\arraystretch}{1.25}
    \begin{tabular}{c c p{0.34\textwidth} p{0.25\textwidth}}
      \textbf{CLA} & \textbf{INS} & \textbf{Commande} & \textbf{Role} \\
      \hline
      80 & 10 & GetBalance(PIN, Solde) & Consultation du solde \\
      80 & 20 & Debit(PIN, Montant, Solde) & Paiement en FCFA \\
      80 & 30 & Credit(PIN, Montant, Solde) & Rechargement \\
    \end{tabular}
  \end{center}
  \vspace{0.4cm}
  \begin{block}{Codes retour}
    9000 : succes \quad
    6B01 : solde insuffisant \quad
    6B02 : montant invalide \quad
    6B03 : PIN incorrect \quad
    6B04 : carte bloquee
  \end{block}
\end{frame}

% --------------------------------------------------------------------
\begin{frame}[fragile]{Securite du PIN}
  \begin{columns}[T]
    \begin{column}{0.50\textwidth}
      \begin{itemize}
        \item Le PIN est verifie dans la carte.
        \item Les erreurs sont comptees en EEPROM.
        \item Apres 3 erreurs : blocage.
        \item Un PIN correct remet le compteur a zero.
      \end{itemize}
    \end{column}
    \begin{column}{0.48\textwidth}
\begin{lstlisting}
If PinSaisi <> PinCode Then
    PinErrors = PinErrors + 1
    If PinErrors >= MaxPinErrors Then
        CarteBloquee = 1
        SW1SW2 = swCardBlocked
        Exit
    End If
    SW1SW2 = swInvalidPIN
    Exit
End If
\end{lstlisting}
    \end{column}
  \end{columns}
\end{frame}

% --------------------------------------------------------------------
\begin{frame}[fragile]{Paiement et ticket}
  \begin{columns}[T]
    \begin{column}{0.48\textwidth}
      \begin{block}{Scenario}
        \begin{itemize}
          \item L'utilisateur saisit un montant.
          \item Le terminal appelle \texttt{Debit}.
          \item La carte retire le montant.
          \item Le terminal imprime le ticket.
        \end{itemize}
      \end{block}
    \end{column}
    \begin{column}{0.50\textwidth}
\begin{lstlisting}
---------- TICKET DE PAIEMENT ----------
Pays      : Senegal
Monnaie   : FCFA
Montant   : 500 FCFA
Statut    : Paiement accepte
Nouveau solde : 9500 FCFA
----------------------------------------
\end{lstlisting}
    \end{column}
  \end{columns}
\end{frame}

% --------------------------------------------------------------------
\begin{frame}{Demonstration}
  \begin{enumerate}
    \item Lancer \texttt{compiler.bat}.
    \item Lancer \texttt{lancer.bat}.
    \item Entrer le PIN \texttt{1234}.
    \item Consulter le solde : \texttt{10000 FCFA}.
    \item Payer \texttt{500 FCFA}.
    \item Verifier le ticket et le nouveau solde.
    \item Tester 3 mauvais PIN pour montrer le blocage.
  \end{enumerate}
\end{frame}

% --------------------------------------------------------------------
\begin{frame}{Resultats obtenus}
  \begin{columns}[T]
    \begin{column}{0.48\textwidth}
      \begin{block}{Compilation}
        \begin{itemize}
          \item \texttt{bin/MaCarte.img}
          \item \texttt{bin/MonTerminal.exe}
          \item \texttt{bin/MonTerminal.img}
        \end{itemize}
      \end{block}
    \end{column}
    \begin{column}{0.48\textwidth}
      \begin{block}{Simulation}
        \begin{itemize}
          \item Carte detectee.
          \item PIN verifie.
          \item Paiement en FCFA execute.
          \item Ticket affiche.
          \item Blocage PIN disponible.
        \end{itemize}
      \end{block}
    \end{column}
  \end{columns}
\end{frame}

% --------------------------------------------------------------------
\begin{frame}{Limites et ameliorations}
  \begin{block}{Limites actuelles}
    \begin{itemize}
      \item Pas de chiffrement des transactions.
      \item Pas d'historique complet des paiements.
      \item Pas d'identite commercant/client stockee.
      \item Pas de connexion bancaire reelle.
    \end{itemize}
  \end{block}
  \begin{block}{Ameliorations possibles}
    Ajouter un code administrateur, un historique EEPROM, le nom du client et le nom du commercant sur le ticket.
  \end{block}
\end{frame}

% --------------------------------------------------------------------
\begin{frame}[plain]
  \begin{tikzpicture}[remember picture,overlay]
    \fill[Ink] (current page.south west) rectangle (current page.north east);
    \node[align=center,text=white] at (current page.center) {
      {\Huge\bfseries Merci}\\[0.35cm]
      {\large Questions ?}\\[0.65cm]
      \tagbox{SenegalGreen}{BasicCard}
      \hspace{0.2cm}
      \tagbox{SenegalYellow}{FCFA}
      \hspace{0.2cm}
      \tagbox{SenegalRed}{Senegal}
    };
  \end{tikzpicture}
\end{frame}

\end{document}
